\documentclass{article}
\usepackage[utf8]{inputenc}
\usepackage{geometry}
\geometry{a4paper, margin=1in}
\usepackage{hyperref}

\title{Intelligence as a System: A New Paradigm}
\author{Albert Jan van Hoek}
\date{\today}

\begin{document}

\maketitle

\section*{Intelligence as a System}

Intelligence isn't something you possess. It's something a system does.

Five Nobel laureates, working in different fields and different decades, all arrived at the same quiet revolution:

\begin{itemize}
    \item \textbf{Herbert Simon} and \textbf{Daniel Kahneman} showed we are bounded, biased, and systematically unreliable as solo thinkers.
    \item \textbf{Friedrich Hayek} proved no single mind can ever hold enough knowledge to act effectively alone.
    \item \textbf{Elinor Ostrom} demonstrated that collective systems only succeed when communication, trust, and real feedback are engineered in.
    \item \textbf{Joseph Stiglitz} revealed that when information is asymmetric or corrupted, even the most capable actors produce systemic failure.
    \item \textbf{Paul Romer} showed that long-run progress is simply accumulated learning—ideas building on ideas over time.
\end{itemize}

Run their findings to their logical conclusion and the picture is unmistakable: \textbf{Intelligence is not a property of individuals. It is the output of a learning loop.}

\section*{The Learning Loop}

Diverse inputs $\rightarrow$ rapid testing $\rightarrow$ honest correction $\rightarrow$ repeat.

The longer and cleaner that loop runs, the smarter the system becomes. A person who seems exceptionally intelligent has simply been running a higher-quality loop for longer, with better inputs and lower-noise feedback. There is nothing mystical about it.

What makes this explosive is that the loop is \textbf{autocatalytic}. Better intelligence spots better inputs, asks sharper questions, and kills bad feedback faster. Small differences in loop quality compound into massive gaps between individuals, teams, organizations, and entire societies.

\section*{The Right Question}

The right question is never ``How smart are the people?'' It is always: \textbf{How well is the loop running?}

That single shift reframes everything—from talent strategy to institutional design to national competitiveness.

\section*{Public Health: A Case Study}

Public health, at its best, has always understood this. It built the purest large-scale learning loop the world has ever seen: sentinel surveillance, real-time modeling, excess-mortality feedback, open data, adversarial review. When any part of that loop was attacked—political capture, suppressed dissent, siloed information—the entire system’s intelligence collapsed exactly as the theory predicts.

\section*{Applying the Architecture}

The same architecture applies to every complex domain: business strategy, scientific research, governance, technology development.

If you want a smarter organization or a smarter society, stop hunting for genius individuals. Build longer, faster, cleaner loops instead:

\begin{itemize}
    \item Radical transparency of data
    \item Low-cost entry for new ideas
    \item High cost for falsified or suppressed information
    \item Institutional memory that survives election cycles and leadership changes
\end{itemize}

Everything else is downstream.

\section*{Conclusion}

The age of the lone genius is over. The age of engineered collective intelligence has begun.

What loop are you running right now—in your team, your company, your industry?

I’d love to hear how you’re designing or defending yours.

\end{document}